%! Sine, Cosine, and Tangent — Unit 6, Lesson 6.3 — Group Activity
\documentclass[10pt]{article}
\usepackage{precalculus-article}
\usepackage{precalculus-boxes}
\newcommand{\TallMath}[1]{$\displaystyle #1\rule[-1.4em]{0pt}{3.2em}$}

\begin{document}

\pageheader{Unit 6, Lesson 6.3}{Group Activity}
\namepartnerperiod

\vspace{0.10in}

\begin{scenariobox}[Scenario]{plumlight}
A point $P$ moves around a unit circle as angle $\theta$ increases from $0$.
The position of $P$ is completely determined by $\theta$: the $x$-coordinate gives $\cos\theta$
and the $y$-coordinate gives $\sin\theta$.
In this activity you will read these values directly from the unit circle and connect them to the
tangent function.
\end{scenariobox}

\vspace{0.12in}

%-----------------------------------------------------------------------
% TIER R
%-----------------------------------------------------------------------
\begin{tcolorbox}[colback=white, colframe=black!40,
    title=\textbf{Tier R --- Remediate: Reading the Unit Circle},
    fonttitle=\bfseries, arc=2mm, left=4mm, right=4mm, top=3mm, bottom=3mm]

\textbf{Directions:} The table below gives the coordinates of $P = (\cos\theta,\, \sin\theta)$ on
the unit circle for several key angles. Some entries are provided to help you --- fill in all
remaining blanks.

\medskip
\begin{center}
\renewcommand{\arraystretch}{2.2}
\begin{tabular}{|c|c|c|c|c|}
\hline
$\theta$ (radians) & $P = (x,\, y)$ & $\cos\theta = x$ & $\sin\theta = y$ & $\tan\theta = y/x$ \\
\hline
$0$      & $(1,\; 0)$                     & $1$  & $0$  & \blank{2cm} \\
$\pi/6$  & $\bigl(\tfrac{\sqrt{3}}{2},\, \tfrac{1}{2}\bigr)$  & \blank{2cm} & \blank{2cm} & \blank{2cm} \\
$\pi/4$  & \blank{3.2cm}                  & \blank{2cm} & \blank{2cm} & \blank{2cm} \\
$\pi/2$  & $(0,\; 1)$                     & \blank{2cm} & \blank{2cm} & \blank{2cm} \\
$\pi$    & \blank{3.2cm}                  & \blank{2cm} & \blank{2cm} & \blank{2cm} \\
$3\pi/2$ & $(0,\; -1)$                    & \blank{2cm} & \blank{2cm} & \blank{2cm} \\
\hline
\end{tabular}
\end{center}

\medskip
\textbf{Reflection:} For which values of $\theta$ in the table is $\tan\theta$ undefined?
Why? \hrulefill

\vspace{0.18in}
\hrulefill
\end{tcolorbox}

\vspace{0.14in}

%-----------------------------------------------------------------------
% TIER A
%-----------------------------------------------------------------------
\begin{tcolorbox}[colback=white, colframe=black!40,
    title=\textbf{Tier A --- On Level: A Point on the Unit Circle and Beyond},
    fonttitle=\bfseries, arc=2mm, left=4mm, right=4mm, top=3mm, bottom=3mm]

\textbf{Part 1.} A point $P = \bigl(-\tfrac{\sqrt{3}}{2},\, \tfrac{1}{2}\bigr)$ is on the unit circle.

\begin{enumerate}[leftmargin=1.8em, itemsep=4pt]
  \item Verify $P$ is on the unit circle. Show that $x^2 + y^2 = 1$. \vspace{0.35in}
  \item State $\sin\theta$, $\cos\theta$, and $\tan\theta$. \vspace{0.35in}
  \item In which quadrant does the terminal ray lie? How do the signs of $x$ and $y$ tell you? \vspace{0.35in}
\end{enumerate}

\textbf{Part 2.} Now suppose a different point $Q = (-\sqrt{3},\, 1)$ is on a circle of \textbf{radius 2}
centered at the origin.

\begin{enumerate}[leftmargin=1.8em, itemsep=4pt, start=4]
  \item Verify $Q$ is on the circle of radius 2. \vspace{0.30in}
  \item Use the general definitions $\sin\theta = y/r$, $\cos\theta = x/r$, $\tan\theta = y/x$
        to compute $\sin\theta$, $\cos\theta$, $\tan\theta$ for the angle $\theta$ whose terminal
        ray passes through $Q$. \vspace{0.35in}
  \item Compare your answers in \#2 and \#5. What do you notice? \vspace{0.25in}
\end{enumerate}
\end{tcolorbox}

\vspace{0.14in}

%-----------------------------------------------------------------------
% TIER E
%-----------------------------------------------------------------------
\begin{tcolorbox}[colback=white, colframe=black!40,
    title=\textbf{Tier E --- Extend: Using the Pythagorean Identity},
    fonttitle=\bfseries, arc=2mm, left=4mm, right=4mm, top=3mm, bottom=3mm]

\begin{enumerate}[leftmargin=1.8em, itemsep=6pt]
  \item You are told that $\sin\theta = \dfrac{3}{5}$ and that $\theta$ is in \textbf{Quadrant II}.
        \begin{enumerate}[leftmargin=1.4em, itemsep=3pt]
          \item[(a)] Use $\sin^2\theta + \cos^2\theta = 1$ to find $\cos\theta$.
                Justify the sign of your answer. \vspace{0.55in}
          \item[(b)] Find $\tan\theta$. \vspace{0.35in}
        \end{enumerate}

  \item Consider the angle $-\theta$ (a clockwise rotation of the same magnitude as $\theta$).
        The point $P$ for $\theta$ is at $(x, y)$; the point for $-\theta$ is at $(x, -y)$.
        \begin{enumerate}[leftmargin=1.4em, itemsep=3pt]
          \item[(a)] Find $\sin(-\theta)$ in terms of $\sin\theta$. \vspace{0.28in}
          \item[(b)] Find $\cos(-\theta)$ in terms of $\cos\theta$. \vspace{0.28in}
          \item[(c)] Describe the geometric symmetry that explains these results. \vspace{0.30in}
        \end{enumerate}
\end{enumerate}
\end{tcolorbox}

\end{document}
